% !PRINTED 394
\weberpage{394}
%<<CONT
von denen die beiden letzten conjugirt imaginär sind, und folglich, da
sie ausser $x, y, \xi, \xi'$ kein gemeinschaftliches Element haben, nur
noch imaginäre Paare enthalten, von denen wir eins, $\xi_1\eta_1$, nebst
dem dazu conjugirten $\xi_1'\eta_1'$ mit aufgeführt haben.

Es genügt demnach, wenn wir die Existenz von einer weiteren reellen
Doppeltangente beweisen können. Denn diese muss dann in 1) vorkommen
und muss in diesem Complexe zu einem reellen Paare gehören.

Jetzt bilden wir noch die beiden conjugirten Complexe:
\[
\begin{array}{rl}
4) & x\xi_1,\ \xi\eta_1 \\
5) & x\xi_1',\ \xi'\eta_1',
\end{array}
\]
in denen $y$ gewiss nicht vorkommt, da sonst $x, \xi, y$ azygetisch wären
(\S.\,96, 3.). Aus demselben Grunde kommt $\xi$ nicht in 5) und $\xi'$
nicht in 4) vor.

Wenn nun \gesperrt{zunächst} die Complexe 4) und 5) syzygetisch sind, so muss
$\xi_1$ in 5) und $\xi_1'$ in 4) vorkommen, und wenn 4) das Paar $p\xi_1'$
enthält, so muss in 5) das Paar $p\xi_1$ vorkommen. Es ist also $p$ mit
seinem conjugirten Elemente identisch, d. h. \gesperrt{reell}.

Wenn \gesperrt{zweitens} die Complexe 4) und 5) azygetisch sind, so kommen
$\xi_1, \xi$ nicht in 5) vor, und es muss $\eta_1$ in 5) enthalten sein.
Wenn $\eta_1$ in 5) mit einem imaginären Elemente $\xi_2'$ gepaart ist, und
wenn in 4) noch ein weiteres imaginäres Paar $\xi_3\eta_3$ vorkommt, so
setzen sich diese Complexe in folgender Weise fort:
\[
\begin{array}{rl}
4) & x\xi_1,\ \xi\eta_1,\ \xi_2\eta_1',\ \xi_3\eta_3,\ \xi_3'\eta_4 \\
5) & x\xi_1',\ \xi'\eta_1',\ \xi_2'\eta_1,\ \xi_3'\eta_3',\ \xi_3\eta_4',
\end{array}
\]
worin auch $\eta_4$ imaginär ist, und durch die accentuirten Buchstaben
immer die conjugirten Elemente zu den unaccentuirten verstanden sind.
Die beiden Complexe enthalten dann nur noch je ein Paar $\xi_5\eta_5$,
$\xi_5'\eta_5'$, die ein gemeinschaftliches Element enthalten müssen. Es
kann aber nicht $\xi_5 = \eta_5'$ sein, weil sonst auch $\xi_5' = \eta_5$, und
mithin beide Paare identisch wären. Also muss $\xi_5 = \xi_5'$ (oder
$\eta_5 = \eta_5'$) sein; d. h. eine dieser beiden Doppeltangenten
ist \gesperrt{reell}, und damit ist unser Satz 2. bewiesen.

3. \gesperrt{Wenn ausser den vier reellen Doppeltangenten $x, y, p, q$,
deren Existenz der Satz 2. behauptet, noch eine weitere vorhanden ist,
so existiren acht reelle Doppeltangenten, die in einem
Steiner'schen Complexe vier Paare bilden.}
